\section{The test system}
\label{sec:model}

\begin{figure*}[t]
\centering
\includegraphics[width=\linewidth]{fig_model.pdf}
\caption{\textbf{The reduced \syn{} model.} The four metabolic modules (ODE) and three gene-expression modules (jump), and what crosses between them. Arrows run between the blocks, not the boxes they touch. Blue arrows run from gene expression to metabolism, and the drains have a gain of order one. The orange arrow is the weak return. Supplementary File S1 lists each departure from the published model.}
\label{fig:model}
\end{figure*}

We demonstrate \IC{} on one test system at two sizes. The \emph{reduced model} is a
published whole-cell model of \syn{} cut down to its best-measured region. It keeps a real
boundary between ODE metabolism and stochastic gene expression. \emph{M0} is a slice of the
reduced model small enough for an exact posterior to be checked against a known truth.
\Cref{tab:sizes} compares the two and lists the results that use each. All the data are synthetic.

\subsection{The reduced model}
\label{sec:reduced}

We test \IC{} on a model reduced from the well-stirred \syn{} model of
Thornburg and coworkers~\cite{thornburg2022fundamental}. The reduction keeps the best-measured region of
the cell and removes about \pn{model.removed} of it (\cref{fig:model}).

\begin{itemize}
  \item \textbf{Metabolism:} \pn{model.rxn} reactions in \pn{model.odemodules} ODE
  modules: glycolysis to lactate, sugar transport and lactate export, nucleotide
  recycling, and a lumped step for tRNA charging.
  \item \textbf{Gene expression:} \pn{model.jumprxn} reactions in
  \pn{model.jumpmodules} jump modules: transcription, transcript decay and
  translation for \pn{model.genes} genes, and membrane insertion for the one
  membrane protein, the glucose transporter PtsG.
  \item \textbf{State:} \pn{model.species} dynamic species and \pn{model.chemostats}
  held species, each named once in a registry with its initial value, prior width
  and source.
\end{itemize}

The two blocks couple in both directions, as in the published model. Enzyme counts
set metabolic rates, and the costs of transcription and translation drain the
nucleotide and charged-tRNA pools. In the other direction, the pools set the
transcription and translation rate constants at each one-minute rebuild. Membrane
protein counts set the cell's area and volume. No rate constant appears in both
blocks. This property comes from the published model, not from our reduction.

The model makes no claim about \syn{} itself. With most of the cell removed, it
cannot reproduce the doubling time, the metabolite levels or the proteome of the
real organism. The code refuses to report a doubling time, so that no result can
be compared with these quantities by mistake. The model's purpose is to give a
real ODE--stochastic boundary from a published whole-cell model, at a size where
inference can be studied.

\subsection{M0, a slice small enough to check}
\label{sec:m0}

An exact posterior on the full reduced model is projected at about \pn{cost.fullscale}
CPU-hours, above our bound of \pn{cost.kone} CPU-hours per posterior
(\cref{sec:posteriors}). A calibration test needs many posteriors, so we also use M0.
M0 keeps all the metabolic reactions, and the same blocks and code. It cuts gene expression to
\pn{mzero.ngenes} of the \pn{model.genes} genes, and it runs for \pn{mzero.horizon}\,s
instead of the reduced model's cycle of \pn{model.cycle}\,s. The two genes are GAPD, whose
promoter is the best determined, and ptsG, the membrane transporter. M0 holds
\pn{mzero.heldenz} enzyme counts, at their measured copy numbers, and \pn{mzero.heldpts}
carrier totals fixed. These holds are declared reductions.

M0 has \pn{mzero.npar} free parameters: the two promoter strengths, the transcript decay rate,
the forward catalytic constants of enolase (ENO) and aldolase (FBA), and the bulk assay noise
$\sigb$. All the posteriors in this paper are on M0. The predictions on the reduced model reuse
an M0 posterior for the parameters the two models share. They draw the promoter strengths of
the other genes from their prior, at three prior widths, and a viability rule drops the draws
that fail (\cref{sec:budget}).

\begin{table}[t]
\caption{\textbf{The two sizes of the test system.} The reduced model's parameters are those
varied in its predictions: the \pn{budget.nI} inferred on M0, and the promoter strengths of
the genes M0 removes, drawn from the prior and filtered for viability.}
\label{tab:sizes}
\centering
\small
\begin{tabular}{@{}P{1.9cm}P{2.6cm}P{2.6cm}@{}}
\toprule
 & Reduced model & M0 \\
\midrule
Metabolic reactions & \pn{model.rxn} & \pn{model.rxn} \\
Stochastic genes & \pn{model.genes} & \pn{mzero.ngenes} \\
Horizon & \pn{model.cycle}\,s & \pn{mzero.horizon}\,s \\
Parameters & \pn{budget.nI} from M0, \pn{budget.nH} from the prior & \pn{mzero.npar}, inferred \\
Results & unscored predictions (\cref{sec:budget,sec:trajectories}) & posteriors (\cref{sec:posteriors,sec:info}), scored predictions (\cref{sec:predm0,sec:trajectories}) \\
\bottomrule
\end{tabular}
\end{table}

\subsection{Synthetic data and the test}
\label{sec:data}

Each synthetic dataset follows \pn{mzero.cells} cells of M0, simulated at a known parameter set
and observed every 60\,s. It has two streams: the exact transcript count of each gene in each
cell, and a bulk panel of \pn{mzero.npanel} metabolite pools, each the mean over cells with
lognormal assay noise (\cref{sec:obs}). The priors are asserted.

Results marked \emph{unscored} are not compared with any truth. For the scored results, the
test is calibration: over many datasets, a 90\% credible interval must contain the truth about
90\% of the time \cite{talts2018validating}.
